
\section{Experiment}
\label{sec:exp}

\subsection{Settings}

\paragraph{Data.} We collect multi-view data from 3 types of cameras, including industry camera (THumanMV~\cite{zhou2024gps}), mobile phone (4K4D~\cite{xu20244k4d} and SelfCap~\cite{xu2024representing}) and our captured GoPro data. 
%
For training, we take 15 training sequences of industry camera data, 6 sequences of GoPro data, and around 3000 frames from 4K4D dance sequence and SelfCap yoga sequence. 
%Totally, we have over 10,000 frames to train our networks.
Compared to public datasets, we capture large scenes accommodating sports movement of multi-person with a portable GoPro system.  
%
To evaluate our Splat-SAP, we take the sequences of unseen characters or of unseen motions from each dataset.
%
In particular, we train only one model of our affinity module for stage 1, while we train one refinement module per camera type.
%
We pick 6 cameras, facing to characters. %from each dataset.
%
The leftmost and the rightmost cameras are fed into the network as source input views.
%
The other 4 views are used as supervision during training and to compute metrics during evaluation.

\paragraph{Metrics.} For rendering, the quality of synthesized images is measured with widely used PSNR, SSIM~\cite{wang2004ssim} and LPIPS~\cite{zhang2018lpips}.
%
We apply Chamfer distance of both directions to evaluate the quality of geometry.
%
We note that the ground truth point set is reconstructed by using Structure-from-Motion~\cite{schoenberger2016sfm} with all 6 views under a long-time optimization.

\paragraph{Implementation Details.} We employ a two-stage training strategy.
%
We first train the affinity learning module for $100k$ iterations with full training data.
%
For each camera type, we further train the rendering module for $60k$ iterations in stage 2.
%
Our networks can be trained on a single RTX 3090 GPU with 24GB.
%
We set $\alpha = 0.8$ in Eq.~\ref{eq:color}, $\beta_1 = 0.8,\ \beta_2 = 0.2$ in Eq.~\ref{eq:render_loss}, $\gamma = 0.5$ in Eq.~\ref{eq:loss1}, and $\lambda_1 = 0.5,\ \lambda_2 = 0.5$ in Eq.~\ref{eq:loss2}.
%
The input to the first stage is in coarse resolution of $W = 512, H=288$, while the fine resolution of $W=1024, H=576$ for the second stage. 
%
We render the high-resolution image of $W=1280, H=720$ in the end.  



\subsection{Results}

\begin{figure}[htpb!]
  \centering
  \includegraphics[width=1.\linewidth]{figs/4_seq_res.pdf}
  \vspace{-8mm}
  \caption{\textbf{Qualitative comparison of rendering on a sequence of data.} Our method preserves temporal and view consistency against 4D-GS and NoPoSplat.}
  \label{fig:compare_noposplat}
  \vspace{-6mm}
\end{figure}


\paragraph{Baselines.} We compare Splat-SAP with state-of-the-art methods of feed-forward rendering, including NeRF-based method ENeRF~\cite{lin2022enerf}, as well as Gaussian-based methods MVSplat~\cite{chen2024mvsplat}, MVSGaussian~\cite{liu2024mvsgaussian} and NoPoSplat~\cite{ye2024nopose}.
In addition, we compare with the optimization-based method 4D-GS~\cite{Wu2024_4dgs}, which requires a long time optimization on sequential data.
We train ENeRF, MVSplat and MVSGaussian from scratch with the same data setting as our second stage training.
We take the pretrained checkpoint of NoPoSplat provided by the original authors and fine-tune it with our training data.
We feed the fine-resolution inputs to ENeRF and MVSGaussian, while the coarse ones to MVSplat and NoPoSplat due to the high memory cost.

For geometry, we compare with scale-invariant methods such as DUSt3R~\cite{wang2024dust3r}, MASt3R~\cite{leroy2024grounding} and VGGT~\cite{wang2025vggt}. 
Furthermore, the scale-aware method Pow3R~\cite{jang2025pow3r} and metric depth estimation method Prompt-DA~\cite{lin2025prompting} are considered for comparison.
Similar to us, Pow3R requires camera calibration as an auxiliary input.
Particularly, Prompt-DA feeds images along with corresponding coarse depth maps to the network as inputs.

\paragraph{Rendering Comparisons.} We report the quantitative results on datasets of Camera, GoPro and Mobile Phone in Tab.~\ref{tab:num_compare_bg}.
%
Our method, in general, outperforms others, especially on Camera and GoPro datasets.
%
Since LPIPS~\cite{zhang2018lpips} is sensitive to higher resolution, our rendering in the resolution of $W=1280$ is on par with the results of MVSGaussian in the resolution of $W=1024$ on Camera data.
%
However, MVSGaussian and ENeRF can not handle thin structures and result in some missing parts in Fig.~\ref{fig:quantative_comp}(a,b), due to the large sparsity.
%
For mobile data, mobile phones are in the mode of alternate zoom-in and zoom-out, which increases the difficulty for feed-forward methods. 
%
Although the setting is tough, our method still renders the fine-grained results with respect to MVSplat in Fig.~\ref{fig:quantative_comp}.
%
Due to the lack of geometry regularization term in Eq.~\ref{eq:geo_regular}, two pieces of Gaussians are sometimes mis-aligned for MVSplat and NoPoSplat in Fig.~\ref{fig:quantative_comp}(c) and \ref{fig:compare_noposplat}(b). 
Thanks to the geometry foundation model MASt3R and our coarse-to-fine learning strategy, we manage the case of sparse-view camera inputs and preserve the view consistency.


NoPoSplat~\cite{ye2024nopose} also leverages MASt3R as a geometry prior, but it gets rid of the traditional stereo constraint, leading to a bad perspective on target view in Fig.~\ref{fig:compare_noposplat}(b). 
During the inference, NoPoSplat still requires source and target view camera poses and normalizes them into a relative scale. 
In Fig.~\ref{fig:compare_noposplat}(b), such normalization leads to rendering jitters in the case of dynamic differences caused by human movement.
%
4D-GS~\cite{Wu2024_4dgs} also neglects geometric constraint and struggles to achieve temporal consistency, see Fig.~\ref{fig:compare_noposplat}(a), for fast motion under sparse views, even if it optimizes on the sequential data for a long time.

\begin{figure}[t!]
  \centering
  \includegraphics[width=1.\linewidth]{figs/4_abla_geo.pdf}
  \vspace{-6mm}
  \caption{\textbf{Qualitative comparison of geometry.} We show point maps of (a) DUSt3R, (b) VGGT, (c) Ours without pixel-wise translation, and (d) Our full affinity. Here is the point map reconstruction with corresponding pixels.}
  \label{fig:geo_ablation}
  \vspace{-3mm}
\end{figure}


\paragraph{Geometry Comparisons.}

As scale-invariant methods, the original point maps of DUSt3R, MASt3R and VGGT are defined in canonical space.
%
Therefore, we employ the ground truth scale factor by comparing their bounding box with that of the ground truth to rescale them into real space.
%
However, DUSt3R sometimes immerses into a local minimum, and projects foreground points onto background, see Fig.~\ref{fig:geo_ablation}(a). 
%
VGGT is not able to handle two-view input with large sparsity, which leads to misalignment of the foreground human from two input views in Fig.~\ref{fig:geo_ablation}(b).  
%
Further, the scale of the scene can not be perfectly estimated by Pow3R, even using camera calibration, thus causing a large Chamfer distance in Tab.~\ref{tab:geometry}. 
%
In addition, we feed the rescaled result of MASt3R as coarse depth input to Prompt-DA.
But the diffusion-based method increases the uncertainty of prediction, and can not preserve 3D consistency from two input views.
Although our method is trained without any geometry loss, we still achieve a superior result in Tab.~\ref{tab:geometry}.
We notice that the ground truth of geometry is a relatively sparse point cloud when using SfM under 6 input views, thus the Chamfer distance from ground truth to prediction can better reflect geometry quality.
%

\begin{figure}[htpb!]
  \centering
  \includegraphics[width=1.\linewidth]{figs/4_abla_full.pdf}
  \vspace{-4mm}
  \caption{\textbf{Qualitative ablation results.} 
  \textit{Upper row:} the rendering comparison between (a) stage 1 and (b) stage 2.
  \textit{Bottom row} illustrates the effectiveness of our depth refinement module in stage 2: (c) initial depth map rendered by affine transformed point maps, (d) depth map after refinement, and (e) rendering results.}
  \label{fig:abla_full}
  \vspace{-2mm}
\end{figure}




\begin{table}[t!]
\small
\centering
\begin{tabular}{l|cc}
\toprule

\multicolumn{1}{c}{\multirow{1}[0]{*}{Method}} & Pred $\rightarrow$ GT $\downarrow$ & GT $\rightarrow$ Pred $\downarrow$ \\ \midrule

DUSt3R         & 0.305 & 0.160 \\
VGGT           & 0.288 & 0.129 \\
Pow3R          & 0.281 & 0.134 \\
MASt3R         & 0.212 & 0.069 \\
Prompt-DA      & 0.205 & 0.063 \\
\hline 
Ours \textit{w/o} Translation   & 0.191 &  0.046 \\
Our Full Model   & \textbf{0.172}  & \textbf{0.027}  \\

\bottomrule
\end{tabular}
\vspace{-2mm}
\caption{\textbf{Quantitative comparison of geometry.} For scale-invariant methods, we compute the rescale factor by comparing their bounding box with that of the ground truth.}
\vspace{-4mm}
\label{tab:geometry}
\end{table}


\paragraph{Ablation Study.} 
We first evaluate the effectiveness of our pixel-wise translation in stage 1.
Our point maps are rescaled with 3-dimensional scaling factors (Eq.~\ref{eq:scale}) when considering camera intrinsic embedding.
%
The point maps can still not avoid misalignment by only using the scaling operator, see Fig.~\ref{fig:geo_ablation}(c).
%
Integrating iterative pixel-wise translation learning, our full affinity module further improves the Chamfer distance in Tab.~\ref{tab:geometry}, and corrects misaligned parts in Fig.~\ref{fig:geo_ablation}(d). 
%


In addition, we evaluate the effectiveness of our refinement module in stage 2. 
%
Alternatively, we can directly synthesize the target view with two Gaussian planes learned by auxiliary layers in stage 1.
%
However, such models typically struggle with holes in the boundary area between foreground and background, see Fig.~\ref{fig:abla_full}(a).
%
The refinement module can correct artifacts and refine details, see Fig.~\ref{fig:abla_full}(b, d).
